\documentclass[11pt]{article}
\usepackage[margin=1in]{geometry}
\usepackage{amsmath,amssymb,graphicx,xcolor}
\usepackage[hidelinks]{hyperref}
\usepackage{caption}
\captionsetup{font=small,labelfont=bf}
\newcommand{\az}{a_0}
\newcommand{\Ez}{E(z)}
\newcommand{\msun}{M_\odot}

\title{\bf Does the MOND scale track the apparent or the event horizon?\\
An evolving-$a_0$ framework, its one theory advance,\\ and why current data cannot yet decide}
\author{Carl Zimmerman\\[2pt]\small\it with computational review support}
\date{June 2026}

\begin{document}
\maketitle

\begin{abstract}
\noindent
The Milgromian acceleration $\az\simeq1.2\times10^{-10}\,\mathrm{m\,s^{-2}}$ coincides with $cH_0$ to
order unity. We take this seriously as a statement that $\az$ is the surface gravity of the cosmic
horizon, $\az=\tfrac{c}{2}\sqrt{G\rho}=cH/Z$ with $Z=2\sqrt{8\pi/3}$, and ask the one question that
makes it falsifiable: does $\az$ track the \emph{apparent} (Hubble) horizon, so that $\az(z)=\az(0)\Ez$
\emph{rises} with redshift, or the \emph{event} (de Sitter) horizon, so that $\az\propto\sqrt{\Lambda}$
is \emph{constant}? We report an honest accounting. On the theory side we make one concrete advance: the
deep-MOND \emph{sign} requires a flat density of states, which is supplied by the $(1{+}1)$-dimensional
near-horizon (s-wave) sector and not by the $3$D bulk; we identify the spectrum and a candidate
probe-coupling selection rule, but stop short of a derivation, which reduces to a defined calculation in
the double-scaled SYK dual. On the data side, the External Field Effect---the one ${\Lambda}$CDM-hard
MOND signature---is present in SPARC at ${\sim}4.8\sigma$ and survives adversarial tests. The direct
$\az(z)$ test, however, is \emph{inconclusive}: re-analysing the two largest real high-$z$ kinematic
samples (KMOS$^{\rm 3D}$, KROSS) we find the curated ``rise'' does not survive, the central estimates
lean toward constant $\az$, but an irreducible gas$+V_{\rm flat}$ systematic spans the entire
framework-versus-constant gap. We conclude that the distinctive prediction is neither supported nor
excluded by existing data, and that the decisive test is a dedicated deep-MOND-tail measurement at
$z\simeq1$--$3$ with a per-galaxy molecular$+$atomic gas census. We state the framework's open core and
its weaknesses explicitly, and we do not connect $\az$ to the Standard Model.
\end{abstract}

\section{The coincidence, taken literally}
That the MOND scale satisfies $\az\sim cH_0$ is an old observation \cite{milgrom1983,famaey2012}. In
emergent-gravity pictures \cite{jacobson1995,verlinde2011} an infrared acceleration scale tied to the
cosmic horizon is natural rather than accidental. We write the scale as the surface gravity of the
free-fall horizon at the radius $R_*=c/\sqrt{G\rho}$ set by the ambient density $\rho$,
\begin{equation}
  \az=\frac{c}{2}\sqrt{G\rho}=\frac{c^2}{2R_*}=\frac{cH}{Z},\qquad Z=2\sqrt{8\pi/3}\simeq5.79,
  \label{eq:a0}
\end{equation}
where the second equality uses $\rho=\rho_{\rm crit}=3H^2/8\pi G$. With the observed $\az$ this gives
$H_0\simeq71\,\mathrm{km\,s^{-1}Mpc^{-1}}$, in the late-time camp but not competitive with the distance
ladder once the unpinned factor $Z$ (the surface-gravity ``$2$'') and the ${\sim}20\%$ stellar $M/L$
systematic are included. Equation~(\ref{eq:a0}) reproduces the standard MOND phenomenology---the
baryonic Tully--Fisher relation $V^4=G M_{\rm b}\az$, the radial-acceleration relation
\cite{mcgaugh2016}, the Freeman surface density $\az/2\pi G$---and we claim nothing there that MOND does
not already claim. Our single new physical content is the \emph{evolution} implied by tying $\az$ to the
\emph{instantaneous} cosmic density.

\paragraph{The one falsifiable question.} Equation~(\ref{eq:a0}) with $\rho=\rho_{\rm crit}(z)$ gives
$\az(z)=\az(0)\Ez$ with $\Ez=\sqrt{\Omega_m(1+z)^3+\Omega_\Lambda}$: $\az$ \emph{rises}. But the standard
emergent-gravity reading ties $\az$ to the cosmological constant, $\az\sim c\sqrt{\Lambda/3}$, which is
\emph{constant}. These are the \emph{apparent} (Hubble) and \emph{event} (de Sitter) horizons
respectively; they agree at $z=0$ to within $\sqrt{\Omega_\Lambda}=0.83$, smaller than the systematic
floor, and diverge by $E(3)\simeq4.6$ by $z\simeq3$. The entire framework reduces to: \emph{which horizon
sets the infrared scale of gravity?} Everything below is an attempt to answer it honestly.

\section{The deep-MOND sign: identified, not derived}
The hardest open problem in any emergent-MOND program is the \emph{sign}: gravity must be
\emph{enhanced} at low acceleration. A temperature-modified Clausius relation gives the wrong sign
(anti-MOND); the right sign comes from \emph{freezing degrees of freedom}. On a holographic screen with
$N=Ac^3/G\hbar$ bits and equipartition $Mc^2=\tfrac12 N_{\rm eff}k_BT$, freezing out a fraction
$N_{\rm eff}/N\to(a/\az)$ below $\az$ raises the temperature required to store the energy and yields
$a=\sqrt{\az g_N}$---enhanced, with $\az\sim cH$ once the freezing scale equals the de Sitter
temperature $T_{\rm dS}=\hbar H/2\pi k_B$.

\paragraph{What forces the freezing law.} The freezing is linear, $N_{\rm eff}\propto T$, only for a
\emph{flat} density of states $g(E)=\mathrm{const}$: writing $g(E)\propto E^p$ gives
$N_{\rm eff}(T)\propto T^{p+1}$, and deep MOND requires $p=0$. We computed the relevant spectra in the
de Sitter static patch by the 't~Hooft brick-wall method \cite{thooft1985}. The \emph{full $3$D} bulk
modes give $g(E)\propto E^2$ ($p=2$, $N_{\rm eff}\propto T^3$)---the area-law entropy and the correct
$T_{\rm dS}$, but \emph{not} MOND. The \emph{$(1{+}1)$D near-horizon s-wave} gives $g(E)=\mathrm{const}$
($p=0$): the flat spectrum that yields the deep-MOND law. Near-horizon dimensional reduction (the
membrane paradigm; the $2$D conformal anomaly derivation of Hawking flux; the near-horizon CFT whose
Cardy entropy is the area law) makes the relevant degrees of freedom $(1{+}1)$-dimensional.

\paragraph{What we do not claim.} The area-law entropy lives in the $3$D modes while MOND lives in the
s-wave, so the two are not the same set; we offer a resolution---a spherically symmetric probe sources
only the monopole ($\ell=0$), whose sector is the flat one, so the \emph{force} is mediated by the
s-wave while the \emph{entropy} is carried by all $\ell$---but this is a selection-rule argument, not a
proof, and the $\az/cH$ coefficient that a sharp cutoff predicts ($3/\pi^2$, i.e.\ $\az\simeq0.30\,cH$)
is the right order but $\sim$1.7$\times$ the observed value, so $Z$ is \emph{not} pinned. The honest
status is: the sign is \emph{identified} (flat $(1{+}1)$D spectrum, Gibbons--Hawking cutoff) and
\emph{motivated}, but \emph{not derived}. The remaining steps---the equipartition projection onto
$\ell=0$ and the exact cutoff---are a defined calculation in the double-scaled SYK (DSSYK) theory that
is the proposed dual of the $2$D de Sitter static patch \cite{narovlansky2023}; we verified one
necessary condition, that the DSSYK density of states is flat near the de Sitter point.

\section{The theoretical lean is conditional}
Rigorous FRW horizon thermodynamics favours the apparent horizon: applying $\delta Q=T\,dS$ on the
apparent horizon $R_A=c/H$ reproduces the Friedmann equation exactly (Cai--Kim \cite{caikim2005}),
whereas the event horizon is teleological and fails the first and generalized second laws in a general
FRW universe. Thus ``$\az$ from horizon thermodynamics'' selects the apparent horizon, i.e.\ rising
$\az$. But this is a \emph{lean, not a proof}, on two counts. First, Cai--Kim concerns the cosmic
\emph{background} (Friedmann); the link to the \emph{local} $\az$ requires a second ingredient, the
Deser--Levin \cite{deserlevin1997} de Sitter--Unruh floor $T(a)=(\hbar/2\pi ck_B)\sqrt{a^2+(cH)^2}$,
which saturates at $T_{\rm dS}$ and places the MOND knee at $a\sim cH$. Second, the constant-$\az$
reading rests on a different but internally consistent premise---$\Lambda$ as the fundamental infrared
cutoff---which the thermodynamics does not exclude. So theory provides a coherent, two-ingredient
package \emph{for} the rising reading and does not exclude the alternative; it does not decide the
question.

\section{Empirical status}
\subsection{The External Field Effect: a robust leg}
The External Field Effect (EFE) is the one signature $\Lambda$CDM does not naturally produce: because
MOND is nonlinear it violates the Strong Equivalence Principle, so a uniform external field suppresses a
system's internal MOND boost \cite{chae2020}. In the SPARC radial-acceleration relation
\cite{lelli2016} we find the outer-disk residual turning negative ($\sim$$-0.1$~dex) with growing
scatter ($0.16\!\to\!0.26$~dex) below $g_{\rm bar}\sim e_N\az$ with $e_N\sim0.03$, at
$\sim4.8\sigma$ when the significance is deflated to the number of \emph{galaxies}. It survives
adversarial tests: fixing $\az$ on the high-acceleration points alone (so the downturn cannot be
absorbed) leaves it at $3.7$--$6.0\sigma$; it appears in both random halves and survives a jackknife.
The one caveat these cannot remove is the outer-disk kinematic systematic (warps, non-circular motions,
beam smearing), which only a per-galaxy environment correlation settles. We regard the EFE as genuine,
foundation-level evidence that gravity is Milgromian, independent of the evolution question.

\subsection{The direct $\az(z)$ test is inconclusive}
A curated three-point compilation (local SPARC; Varasteanu 2025; a MUSE-DARK $z\!=\!0.9$ point) favours
rising $\az$ at a naive $\sim5\sigma$, falling to $\sim2.5\sigma$ once a redshift-growing gas/pressure
systematic is marginalized, and to a coin toss if the single $z\!=\!0.9$ point is dropped. We therefore
re-analysed the two largest \emph{real} high-$z$ kinematic samples directly.
\begin{itemize}\itemsep2pt
\item \textbf{KMOS$^{\rm 3D}$} (\"Ubler et al.\ 2017 \cite{ubler2017}; 135 galaxies, $z=0.6$--$2.5$):
the baryonic-Tully--Fisher $\az=V_c^4/(G M_{\rm b})$ is ${\sim}4\times$ the local value and \emph{rises
with mass}---a high-acceleration bias---and its redshift trend is flat-to-declining, not rising. The
sample is too massive and too dispersion-supported to measure $\az$ cleanly.
\item \textbf{KROSS} (Harrison et al.\ 2017 \cite{harrison2017}; 381 rotation-dominated galaxies,
$z\simeq0.85$): with a realistic molecular-gas correction \cite{tacconi2018} the BTFR normalization
gives $\az(0.85)$ consistent with the \emph{local} (constant) value, below the framework's predicted
$1.96\times10^{-10}$. Only an unphysical no-gas treatment matches the rise.
\end{itemize}
Crucially, decomposing the irreducible systematics---the CO-to-H$_2$ factor, the \emph{unmeasured}
atomic gas, and whether $V_c$ reaches $V_{\rm flat}$ (which $\az\propto V^4$ amplifies)---yields
$\az(0.85)\in[0.5,2.5]\times10^{-10}$, \emph{wider} than the framework ($1.96$) to constant ($1.20$)
gap. The central estimate leans constant, but the systematic spans both models. The honest verdict is
\emph{inconclusive with a central preference for constant}: the curated ``rise'' does not survive, but
the data do not exclude the framework either. No existing sample resolves this, because the
gas-measured samples are massive/high-surface-brightness (wrong regime) while the kinematic samples lack
per-galaxy gas.

\section{Weaknesses, stated plainly}
(i) \emph{Clusters.} MOND's residual-mass problem (a factor ${\sim}2$ in relaxed clusters) is not
relieved by the cosmic $\az(z)$, which boosts the dynamical mass only as $\sqrt{\Ez}$ ($\sim$8\% at
$z=0.3$) and cannot be made environment-dependent without breaking the tight galaxy RAR.
(ii) \emph{$\az$ universality.} A per-galaxy SPARC analysis rejects a strictly universal $\az$ at high
formal significance ($\chi^2/\mathrm{dof}\!\sim\!770$ statistical), reduced but not erased by the
standard $0.13$~dex systematic floor ($\chi^2/\mathrm{dof}\!\sim\!5$); the residual scatter correlates
with mass at $r\!\sim\!0.37$, suggestive of a density dependence but degenerate with fitting artifacts.
(iii) \emph{Second-order CMB.} Linear $C_\ell$ are $\az$-independent, but rising $\az$ drives
recombination into the deep-MOND regime (whereas constant $\az$ leaves it Newtonian), amplifying a
second-order, non-Gaussian imprint through the non-analytic $\mathcal{Y}^{3/2}$ term---a potential
discriminator and a potential tension that a second-order Boltzmann calculation must settle.

\section{The decisive test}
The clean experiment is a direct $\az$ measurement in a population of rotation-supported, deep-MOND-tail
discs at $z\simeq1$--$3$, analysed as a stellar-$M/L$- and $Z$-cancelling \emph{ratio} $\az(z)/\az(0)$,
with a per-galaxy molecular$+$atomic gas census and resolved rotation curves reaching $V_{\rm flat}$.
The three readings separate as $\az(3)/\az(0)=4.6$ (apparent/rising), $1.0$ (event/constant), $0.42$
(scale-invariant vacuum \cite{maeder2017}, falling). An independent, kinematics-free cross-check is the
weak-lensing prediction that the phantom mass at fixed baryonic mass scales as $\sqrt{\Ez}$
($+34\%$ by $z=1$), accessible to Euclid/Rubin/JWST stacking---immune to the pressure systematic though
not to the gas census. We stress that these are the experiments that decide the question; existing data
do not.

\section{Conclusion}
The surviving, defensible content is one equation, Eq.~(\ref{eq:a0}): the MOND scale as cosmic-horizon
surface gravity, with the standard MOND phenomenology following and one new, falsifiable consequence---an
evolving $\az$. The deep-MOND sign is identified with the flat $(1{+}1)$D near-horizon spectrum and a
candidate coupling rule, but not derived; the derivation reduces to a defined DSSYK calculation. The
apparent-horizon reading is theoretically favoured but only conditionally. Empirically the EFE is a
robust, $\Lambda$CDM-hard detection, while the distinctive evolving-$\az$ prediction is currently
\emph{undecided}: the real high-$z$ data are gas-systematics-limited, with central values leaning
constant. We have deflated the framework's own claims---favourable and unfavourable---wherever they
overreached, and we connect $\az$ to gravity alone and not to the Standard Model. Whether the infrared
scale of gravity tracks the apparent or the event horizon is, on the present evidence, an open and
sharply testable question.

\begin{thebibliography}{99}\small
\bibitem{milgrom1983} M.~Milgrom, \textit{ApJ} \textbf{270}, 365 (1983).
\bibitem{famaey2012} B.~Famaey \& S.~McGaugh, \textit{Living Rev.\ Rel.} \textbf{15}, 10 (2012).
\bibitem{jacobson1995} T.~Jacobson, \textit{PRL} \textbf{75}, 1260 (1995).
\bibitem{verlinde2011} E.~Verlinde, \textit{JHEP} \textbf{04}, 029 (2011).
\bibitem{mcgaugh2016} S.~McGaugh, F.~Lelli \& J.~Schombert, \textit{PRL} \textbf{117}, 201101 (2016).
\bibitem{thooft1985} G.~'t~Hooft, \textit{Nucl.\ Phys.\ B} \textbf{256}, 727 (1985).
\bibitem{narovlansky2023} V.~Narovlansky \& H.~Verlinde, arXiv:2310.16994 (2023).
\bibitem{caikim2005} R.-G.~Cai \& S.~P.~Kim, \textit{JHEP} \textbf{02}, 050 (2005).
\bibitem{deserlevin1997} S.~Deser \& O.~Levin, \textit{Class.\ Quantum Grav.} \textbf{14}, L163 (1997).
\bibitem{chae2020} K.-H.~Chae et al., \textit{ApJ} \textbf{904}, 51 (2020).
\bibitem{lelli2016} F.~Lelli, S.~McGaugh \& J.~Schombert, \textit{AJ} \textbf{152}, 157 (2016).
\bibitem{ubler2017} H.~\"Ubler et al., \textit{ApJ} \textbf{842}, 121 (2017).
\bibitem{harrison2017} C.~Harrison et al., \textit{MNRAS} \textbf{467}, 1965 (2017).
\bibitem{tacconi2018} L.~Tacconi et al., \textit{ApJ} \textbf{853}, 179 (2018).
\bibitem{maeder2017} A.~Maeder, \textit{ApJ} \textbf{834}, 194 (2017).
\end{thebibliography}
\end{document}
